% Authority: exact-scope leaf 19; full PDF page 98; printed p.81. Translation of direct transcription.
% U:p081.u01 | paragraph-start
\par\indent
These two cases are comprised in the following formula, to which we
join a result obtained above:
\begin{equation*}\tag{$\eta$}
\leg{t}{p}=(-1)^{\frac{t-1}{2}}\leg{t}{a},\quad
\leg{u}{p}=(-1)^{\frac{p-1}{4}}.
\end{equation*}
% U:p081.u02 | paragraph-start
\par\indent
The equation $t^2-au^2=p$ gives successively:
\[
t^2\equiv au^2,\quad t^{\frac{p-1}{2}}\equiv a^{\frac{p-1}{4}}u^{\frac{p-1}{2}}\md{p},
\]
whence one sees that $a$ will or will not be a biquadratic residue with respect to
$p$, according as one has $\leg{t}{p}=\leg{u}{p}$ or $\leg{t}{p}=-\leg{u}{p}$.
Putting in place of $\leg{t}{p}$ and $\leg{u}{p}$ the values given by the
expressions $(\eta)$, one can replace the preceding conditions by those that follow:
% U:p081.u03 | labelled-proposition
\par\noindent\hspace*{.12\linewidth}\begin{minipage}{.86\linewidth}
$a$ will or will not be a biquadratic residue with respect to $p$, according as one has:
\end{minipage}
\begin{equation*}\tag{$\theta$}
\leg{t}{a}=(-1)^{\frac{t-1}{2}+\frac{p-1}{4}}\quad\text{or}\quad
\leg{t}{a}=-(-1)^{\frac{t-1}{2}+\frac{p-1}{4}},
\end{equation*}
\noindent
a result that we shall not pause to demonstrate.
% U:p081.u04 | paragraph-start
\par\indent
Let us now put $p=\varphi^2+\psi^2$ (where $\psi$ is supposed even) and equate this
value of $p$ to the one we considered previously. We shall thus have
$p=t^2-au^2=\varphi^2+\psi^2$, whence one concludes by transposing:
\[
t^2-\psi^2=(t+\psi)(t-\psi)=\varphi^2+au^2.
\]
% U:p081.u05 | paragraph-start
\par\indent
There are now two cases to distinguish, according as $\varphi$ is or is not
divisible by $a$. Let us begin with that in which $\varphi$ is not divisible by $a$. Let
$m$ be the greatest common divisor of $\varphi$ and of $u$, which will necessarily be
odd and not divisible by $a$, and put $\varphi=m\varphi'$, $u=mu'$, values whose
substitution in the equation obtained above changes it into this one:
\[
(t+\psi)(t-\psi)=m^2(\varphi'^2+au'^2).
\]
% U:p081.u06 | continuation
This equation shows that $t+\psi$ is composed of two factors, one of which
divides $m^2$, the other $\varphi'^2+au'^2$, and that the same holds for $t-\psi$. If we denote
the factors of $t+\psi$ by $E$ and $K$ and those of $t-\psi$ by $F$ and $L$, we shall have:
\[
\begin{aligned}
t+\psi&=EK, &t-\psi&=FL,\\
m^2&=EF, &\varphi'^2+au'^2&=KL.
\end{aligned}
\]
% U:p081.u07 | continuation
One will ascertain, as in paragraph 4, that the numbers $E$ and $F$ are
relatively prime. It follows from this and from the equation $m^2=EF$ that each of
these numbers is a square, so that:
\[
\leg{E}{a}=1,\quad\leg{F}{a}=1.
\]
